\chapter{The C++26 features maya leans on}
\label{ch:cpp26-features}

\begin{aside}
\maya{} is named ``C++26 for the terminal'' for a reason: a
handful of newer language and library features make the
architecture practical. This chapter is a tour of those features
in isolation: what each adds, how \maya{} uses it, and what the
pre-C++26 alternative would look like.
\end{aside}

\section{Concepts (C++20)}

Already covered in Chapter~\ref{ch:concepts}. Without concepts,
template error messages would be wall-of-text SFINAE diagnostics.
With concepts, the constraint is named; the failure point is
specific.

\maya{} uses concepts for: \code{Program}, \code{ElementLike},
\code{Renderable}, \code{StyleSource}, \code{Subscribable},
\code{Terminal}, \code{Addable}, etc.

\section{Designated initialisers (C++20)}

\begin{lstlisting}
Element panel{
    .v = BoxElement{
        .children = {...},
        .layout = FlexStyle{.direction = Direction::Column},
        .border = Border{.style = BorderStyle::Rounded},
    },
};
\end{lstlisting}

Designated initialisers let you construct aggregates by name,
not by position. The pre-C++20 alternative was a fluent builder
(verbose) or careful brace-init (fragile).

\maya{}'s use: \code{Style}, \code{Color}, \code{Border},
\code{FlexStyle}, \code{Padding} are all aggregates and are
typically constructed with designated init.

\section{Structural NTTPs (C++20, refined C++26)}

The DSL's headline feature. \code{Str}, \code{CTStyle},
\code{BoxCfg} are structural types that can appear as template
parameters. Without them, the DSL would have to encode every
parameter as a separate template specialisation by hand.

\maya{}'s use: every compile-time configuration --- text
literals, styles, padding values, border styles, widths, heights,
flex factors --- is carried in an NTTP.

\section{std::variant (C++17)}

The fundamental sum type. \code{Element}, \code{Cmd<Msg>},
\code{Sub<Msg>}, \code{Event} are all variants. Without them,
the same shape would require an OOP hierarchy plus virtual
dispatch.

\section{std::visit and overload (C++17)}

\code{std::visit} matches a variant against a callable. The
overload pattern (multiple lambdas inheriting into one struct)
makes the visitor readable.

\maya{}'s \code{maya::overload} (in
\file{maya/include/maya/core/overload.hpp}) is the pattern as a
named struct.

\section{std::expected (C++23)}

Value-typed error handling (Chapter~\ref{ch:expected}). Before
C++23, you used Boost.Outcome, the absl variant, or a hand-rolled
expected. The standard version simplifies portability.

\maya{} ships its own \code{Expected} that predates the standard
addition; the API is compatible.

\section{Coroutines (C++20)}

\maya{} does not currently use coroutines, but the
\code{Cmd::task} interface is coroutine-shaped: you could
\code{co\_await} on a task. Future versions may take this on.

\section{std::format (C++20)}

\begin{lstlisting}
auto label = std::format(\"Count: {}\", n);
\end{lstlisting}

A type-safe \code{printf}. Used internally by widget code that
constructs labels and tooltips.

The C++23 \code{std::print} extends \code{format} to direct
stdout output; \maya{} prefers \code{print} where supported.

\section{std::span (C++20)}

A view over a contiguous range, like \code{string\_view} but
generic. Used heavily in \maya{}'s I/O code:

\begin{lstlisting}
Expected<size_t, IoError> read(std::span<std::byte> buf);
\end{lstlisting}

No ownership, no copy. The buffer is supplied by the caller.

\section{Ranges (C++20 and C++23)}

\maya{} uses ranges for view-and-transform pipelines:

\begin{lstlisting}
auto labels = items | std::views::transform([](auto& i) { return i.label; });
\end{lstlisting}

Lazy evaluation: the transform runs as the range is consumed,
not eagerly. Useful for building element trees from large
sources.

C++23 added \code{std::views::join\_with} which the DSL uses to
intersperse separators between list items.

\section{Concepts-aware overload resolution (C++20)}

Two templates with overlapping concepts: the more-specialised
one wins. The DSL relies on this heavily; multiple pipe overloads
target the same operator and the compiler picks by the most-
specific concept.

\section{constinit and consteval (C++20)}

\code{constinit} guarantees a variable is initialised at
compile time; \code{consteval} requires a function to run only
at compile time. The DSL's \code{merge}, \code{parse\_hex},
\code{compute\_box\_cfg} are \code{consteval}; the result is
that runtime code never sees the cost.

\section{Lambda extensions (C++23, C++26)}

C++23 added attributes on lambdas; C++26 is exploring
\code{static} lambdas (no captures, no per-call storage).
\maya{}'s tight loops use lambdas; the future-version
optimisations will improve performance without source changes.

\section{Class types in NTTPs (C++20)}

The key C++20 relaxation that made compile-time strings
possible. A structural type with all-public members can be an
NTTP. Without this, \maya{}'s entire DSL would be impossible.

\section{Class template argument deduction (C++17, refined C++20)}

\begin{lstlisting}
auto e = Element{TextElement{\"hi\"}};            // deduces Element type
\end{lstlisting}

The compiler infers the template arguments from constructor
arguments. \maya{} uses CTAD throughout to keep the syntax
clean.

\section{Three-way comparison operator (C++20)}

\begin{lstlisting}
auto cmp(const Style& a, const Style& b) = default;  // <=>
\end{lstlisting}

A single defaulted spaceship operator generates all six
comparison operators. \maya{}'s structural types use this for
NTTP equality.

\section{std::array (C++11) and constexpr containers (C++20)}

The Unicode width table is a \code{constexpr std::array}.
Lookup is binary search. The whole table sits in read-only data;
no startup cost.

\section{The if-with-init form (C++17)}

\begin{lstlisting}
if (auto r = parse(input); r) {
    use(r.value());
}
\end{lstlisting}

The init clause scopes the variable to the \code{if}/\code{else}
chain. Used in every error-handling site.

\section{P2300 senders/receivers (C++26)}

A long-awaited concurrency model. \maya{} does not depend on it
yet but the \code{Cmd::task} abstraction is sender-shaped; future
versions may integrate.

\section{Reflection (C++26)}

The reflection proposal (P2996) would let \maya{} generate
visitors and serialisers automatically. The current code uses
templates and constexpr to approximate; reflection would simplify.

\section{What it costs to use C++26}

The features above require GCC 14+, Clang 17+, or MSVC v19.40+.
Older compilers cannot build \maya{}. The trade-off:

\begin{itemize}[leftmargin=*]
  \item Library is shorter (fewer workarounds).
  \item Code is easier to read (named constraints, designated
    init).
  \item Distribution is harder (newer compiler required).
\end{itemize}

For a library targeting modern terminals on modern OSes,
requiring a modern compiler is acceptable.

\section{What the pre-C++26 version would look like}

If we backported \maya{} to C++17:

\begin{itemize}[leftmargin=*]
  \item Concepts become SFINAE. Error messages explode.
  \item NTTP-based DSL becomes type-tag-based with
    template specialisations per literal.
  \item Designated init becomes builders or positional init.
  \item Variants exist but visit-with-overload is uglier.
\end{itemize}

The result would be 2--3 times as much code for the same
functionality. C++26 is the right floor.

\section*{Workshop}

\begin{enumerate}[leftmargin=*]
  \item Read \file{maya/include/maya/core/concepts.hpp} and
    list every concept by feature it depends on.
  \item Find a single function in \maya{}'s source that uses
    designated initialisers, NTTPs, ranges, and visit in one
    expression. (There are many candidates.)
  \item Try to compile \maya{} with \code{-std=c++17}. Count the
    errors. (They will be many.)
\end{enumerate}

\begin{recap}
\maya{} leans on concepts, designated initialisers, structural
NTTPs, variants, visit, expected, format, span, ranges, and
\code{consteval}. Each of these features removes a workaround
that pre-C++26 codebases live with. The cost is a modern
compiler requirement; the payoff is a small, readable library.
\end{recap}
